% Generated table; it uses the manuscript preamble's macros (\hltint, \hcell, \altrow, PanelBg).
\begin{xltabular}{\textwidth}{>{\raggedright\arraybackslash}X r*{3}{>{\centering\arraybackslash}p{2.7cm}}}
\caption{\hltint{HlM}{Tests of independence from irrelevant alternatives in the multinomial logit.}}\label{tab:a_iia}\\
\toprule
\headerrow
\hcell{Outcome dropped} & \hcell{Riders kept} & \hcell{Hausman and McFadden $\chi^2$} & \hcell{Small and Hsiao rejections (\%)} & \hcell{Rejections under independence (\%)} \\
\midrule
\endfirsthead
\multicolumn{5}{l}{\cellcolor{HeaderGold}\textcolor{HeaderText}{\textbf{\tablename\ \thetable{} \textit{(continued)}}}} \\
\toprule
\headerrow
\hcell{Outcome dropped} & \hcell{Riders kept} & \hcell{Hausman and McFadden $\chi^2$} & \hcell{Small and Hsiao rejections (\%)} & \hcell{Rejections under independence (\%)} \\
\midrule
\endhead
\midrule
\multicolumn{5}{r}{\textit{Continued on next page}} \\
\endfoot
\bottomrule
\multicolumn{5}{p{0.97\linewidth}}{\footnotesize \textit{Note:} Both tests have 36 degrees of freedom for every outcome dropped. A negative Hausman and McFadden statistic leaves no $p$-value defined, so none is given. The rejection columns give the percentage of random sample splits rejecting at the 5\% level, with the number of usable splits in parentheses, on the observed data and on data simulated under independence.}\\
\endlastfoot
O & 445 & $-$33.89 & 76.0 (25) & 35.0 (20) \\
\altrow
BC & 159 & $-$13.09 & 90.5 (21) & 75.0 (8) \\
KA & 440 & $-$1.29 & 52.0 (25) & 40.0 (20) \\
\end{xltabular}
